
Practitioners routinely avoid horizontal flip augmentation when training models on MNIST digits---Kaggle winning solutions exclude it, PyTorch official tutorials omit it---yet no formal validation explains why or quantifies the harm. We rigorously test this folklore through four controlled hypotheses, evaluating horizontal flip at probabilities {0.3, 0.5, 0.9} against baseline and rotation $\pm 15$° control on MNIST digit classification. Asymmetric digits {2, 3, 5, 6, 7, 9} exhibit statistically significant degradation of 0.72-1.00 percentage points at moderate flip rate (p=0.5), with dose-dependent degradation ranging from 0.37 pp at low flip rate (p=0.3) to 4.10 pp at high flip rate (p=0.9), while symmetric digits {0, 1, 8} remain largely stable at moderate flip rates (<0.2\% change at p=0.5). We observe a perfect dose-response relationship in one experiment (Spearman $\rho$ = -1.0, p < 0.001) and near-perfect in another ($\rho$ = -0.969, p = 1.$97\times 10^{-12})$, indicating a deterministic label noise mechanism: flipped asymmetric digits retain original labels despite visual non-canonicality, creating training examples that degrade test accuracy on canonical digits. Rotation $\pm 15$° augmentation produces no differential effect across three independent validations, confirming semantic invalidity---not general augmentation---as the causal factor. These findings demonstrate semantic validity testing on MNIST and propose a generalizable framework: augmentations must preserve domain-specific class identity to avoid introducing label noise.

